\begin{abstract}
Predicting a crystal structure from its composition requires finding both the lattice and the positions of the atoms within it. Diffusion and flow models can predict both, but they typically refine each candidate over many sampling steps. We introduce CrytalDrift, a model that generates a periodic structure in a single network evaluation. During training, a drifting field draws generated crystals toward reference structures and pushes them away from other generated samples. We build this field using periodic descriptors of atomic environments and show that its value is unchanged when any crystal is rotated, its origin is shifted, or its lattice is written in an equivalent basis. We use the field to train generators that update individual atomic sites or move symmetry-related sites together within occupied Wyckoff orbits. On MP-20, one-step sampling recovered more reference structures than a compact multistep flow when both were given the same short time per composition.
\end{abstract}

\section{Introduction}
Crystal structure prediction starts with a chemical composition and asks how its atoms are arranged in a periodic solid. Several structures can share the same formula, so a practical use of a generative model is to propose a set of candidates for each composition. DiffCSP~\cite{diffcsp} generates candidates by jointly denoising the lattice and the fractional atomic coordinates. FlowMM~\cite{flowmm} uses flow matching to evolve a starting configuration into a crystal. Both methods account for periodicity. Sampling many candidates, however, requires many evaluations of the denoising network or the velocity field. One-step predictors such as uFlowCSP~\cite{uflow} show that a single network evaluation can already recover a substantial fraction of known structures.

Drifting models offer a different way to train a one-step generator. During training, an attractive field draws generated samples toward the data, and a repulsive field pushes them away from other samples produced by the model~\cite{deng}. At inference, the trained network maps a new input directly to a sample. The choice of distance used to construct the field is especially important for atomic structures. Two copies of the same molecule can appear far apart in Cartesian coordinates if one is rotated. SymDrift~\cite{symdrift} showed that an equivariant generator can still receive a different training signal when the reference molecules are independently rotated. The authors addressed this by aligning pairs of structures or by computing the field in an invariant embedding space.

Periodicity introduces a further ambiguity. Consider a lattice with basis vectors \(a\), \(b\), and \(c\). Replacing \(b\) with \(a+b\) gives another basis for the same lattice. The fractional coordinates of the atoms change accordingly, even though every atom remains in the same physical position. A field computed directly from lattice entries and fractional coordinates can therefore assign different weights to two descriptions of the same reference crystal. Space-group symmetry plays a different role. Once a space group and its occupied Wyckoff positions have been chosen, moving one atom constrains where its symmetry partners can go. DiffCSP++~\cite{diffcspplus} uses these constraints during diffusion, while the Wyckoff Transformer~\cite{wyformer} learns a distribution over discrete patterns of occupied sites.

In CrytalDrift, we construct the drifting field from descriptors of periodic crystal structures. We compare descriptors based on interatomic distances and angles with extensions that capture contacts between occupied orbits or correlations in reciprocal space. We show that the field is unchanged when the generated crystal and each reference crystal are independently rotated, shifted, or written in an equivalent lattice basis. We use this field to train two generators. One updates individual atomic sites and the lattice from a random starting structure. The other infers an occupied Wyckoff pattern from the composition and updates the free coordinates of each orbit together with the lattice.

Theorem~\ref{thm:field} establishes invariance of the training field. Theorem~\ref{thm:decoder} shows which symmetries the orbit generator preserves in its output.

We evaluate CrytalDrift on MP-20 under two sampling budgets. Fixing the number of proposals measures the chance of recovering the reference structure in a given number of attempts. Fixing the time per composition also accounts for the cost of each attempt. Our compact multistep flow recovers more targets than the one-step site generator at equal proposal counts, but the site generator recovers more under a short time limit. For the Wyckoff-based generator, we also compare each prediction with its own starting structure. This lets us measure how much the learned update contributes beyond what the starting structures already recover.
